\newif\ifglobal

%%%%%%%%%%%%%%%%%%%%%%%
%%% Compile to view %%%
%%%%%%%%%%%%%%%%%%%%%%%

%%%%%%%%%%%%%%%%%%%%%%%%%%%%%%%%%%%%%%%%%%%%%%%%%%%
%%% Readme:                                     %%%
%%% https://www.overleaf.com/read/bwdjvfjksvjy/ %%%
%%%%%%%%%%%%%%%%%%%%%%%%%%%%%%%%%%%%%%%%%%%%%%%%%%%

% >>> M?: marks a module setting number or important notice

% >>> M4: Set {./relative-path} to external files for this module <<<
\def \modulefiles {./28-LPM}
% To add an external file, use:
% (Replace '---' with actual filename)
% '\input{\modulefiles/tables/---.tex}' for tables
% '\includegraphics{\modulefiles/figures/---}' for figures
% '\subfile{\modulefiles/---}' for register files

% >>> M1: This module is in global project? <<<
% 'YES' - uncomment; 'NO' - comment out
\globaltrue

\ifglobal
    \documentclass[main\_\_CN.tex]{subfiles}

\else
    % Font "FandolSong-Regular" used by ctex does not have GSUB table.
    % It causes warnings that do not affect compilation and can be ignored.
    \PassOptionsToPackage{quiet}{fontspec}

    \documentclass[openany, 10pt]{book}

    % >>> M2: In ./00-shared/preamble-trm-module, also find
    %         settings for visibility of todo notes, watermark <<<
    \usepackage{preamble-trm-module}


    % Enable Chinese typesetting
    \usepackage{ctex}

    % >>> M3: Temporary labels in Chinese
    %         you can add your own labels there <<<
    \myexternaldocument{./00-shared/temp-labels-cn}

\fi %\ifglobal


\setcounter{secnumdepth}{4}

% >>> M5: If this module needs any updates in
% './00-shared/preamble-trm-module.sty'
% (include additional package, etc.), contact your module PM
% or create the file './00-shared/changelog.tex' make the
% necessary updates yourself and document them in this file <<<


\begin{document}


% Sets language into which pre-defined bits of text are translated
\selectlanguage{chinese}

% Do not delete this block
\ifglobal
% Add the following front matter if
% this module is outside of global project
\else
    \InputIfFileExists{readme}{}{}
    {\let\clearpage\relax \listoftodos}
    \tableofcontents
    %\listoftables
    %\listoffigures
\fi


%%%%%%%%%%%%%%%%%%%%%%%%%
%%% TRM Module Begins %%%
%%%%%%%%%%%%%%%%%%%%%%%%%

\hypertarget{lowpowm}{}
\chapter{低功耗管理 (RTC\_CNTL)}\label{mod:lowpowm}

\section{概述}

\chipname{} 拥有一个先进的电源管理单元，可在低功耗协处理器的配合下，灵活打开或关闭芯片的不同电源域，协助客户在芯片工作性能、功耗控制和唤醒延迟之间实现最佳平衡。为了便利用户的使用，\chipname{} 定义了四种最常见的电源域设置组合，对应四种预设功耗模式，可满足用户的常见场景需求，但也同时支持用户对某个电源域的独立控制，以满足一些复杂场景的功耗需求。\chipname{} 内置 2 个超低功耗协处理器，允许芯片在绝大多数电源域均关闭的情况下正常工作，因此可以实现超低功耗。

\section{主要特性}
\chipname{} 的低功耗管理有如下特性：
\begin{itemize}
\item 4 种预设功耗模式，可满足多种典型应用场景需求
\item 高达 16 KB 保留内存 (retention memory)，含快速内存和慢速内存
\item 8 个 32 位保留寄存器 (retention register)
\item 支持 RTC Boot 功能，用于快速唤醒
\item 各功耗模式下均支持超低功耗协处理器
\end{itemize}

在本章节中，我们将首先介绍 \chipname{} 低功耗管理的工作过程，其次介绍芯片的预设低功耗工作模式，最后介绍芯片的 RTC Boot 过程。

\section{功能描述}

\chipname{} 的低功耗管理主要由以下模块实现：

\begin{itemize}
\item 功耗管理单元 (PMU)：控制向模拟、RTC 和数字类电源域的供电。具体细分电源域列表，请见第 \ref{sec:lpm-power-domain} 小节；
\item 电源隔离单元：保证各电源域的独立工作，防止掉电电源域影响其他电源域的工作；
\item 低功耗时钟：为低功耗模式下工作的电源域提供时钟信号；
\item 定时器：
    \begin{itemize}
    \item RTC 定时器：在专用寄存器中记录 RTC 主状态机的状态；
    \item ULP 定时器：在预设时间唤醒超低功耗协处理器。更多详情，请见章节 \ref{mod:ulp} \textit{\nameref{mod:ulp}}。
    \item 触摸传感器定时器：在预设时间唤醒触摸传感器。更多详情，请见章节 \ref{mod:sensor}\textit{\nameref{mod:sensor}}。
    \end{itemize}
\item 8 个 32 位 ``always-on" 保留寄存器：即这 8 个寄存器永远处于工作状态，不受 deep-sleep 等低功耗模式的影响，可用于存储一些不能丢失的数据。

\item 22 个 ``always-on" 管脚：即这 22 个管脚永远处于工作状态，不受 deep-sleep 等低功耗模式的影响，可用作低功耗模式下的唤醒源（详见第 \ref{sec:lpm-wakeup-source} 节），也作为正常 GPIO 使用（详见 \ref{mod:iomuxgpio} \textit{\nameref{mod:iomuxgpio}} 章节）。
\item RTC 慢速内存：8 KB SRAM，工作时钟为 RTC 快速时钟 (rtc\_fast\_clk)，可用作扩展内存或存储 ULP 指令和数据内存。
\item RTC 快速内存：8 KB SRAM，工作时钟与 CPU 时钟 (CPU\_CLK) 同频，可用作扩展内存。
\item 调压器：调节向不同电源域的供电电压。
\end{itemize}

\chipname{} 低功耗管理的原理图可见图 \ref{fig:lpm-esp32-power-control}。
\begin{figure}[H]
\centering
\includegraphics[width=1.0\textwidth]{\modulefiles/figures/28-LPM-power_ctrl.pdf}
\caption{低功耗管理原理图}
\label{fig:lpm-esp32-power-control}
\end{figure}

\textbf{说明：}
\begin{itemize}
    \item 有关各电源域的具体描述，请见第 \ref{sec:lpm-power-domain} 节。
    \item 上图中的信号描述如下：
    \begin{itemize}
        \item xpd\_rtc\_reg：
        \begin{itemize}
            \item 当 \hyperref[fielddesc:RTCCNTLRTCREGULATORFORCEPU]{RTC\_CNTL\_RTC\_REGULATOR\_FORCE\_PU} 置 1 时，低功耗调压器常开；
            \item 否则，低功耗调压器在芯片进入 Light-sleep 和 Deep-sleep 时关闭。此时，RTC 电路由一个极低功耗的内置电源供电。
        \end{itemize}
        \item xpd\_dig\_reg：
        \begin{itemize}
            \item 当 \hyperref[fielddesc:RTCCNTLDGWRAPPDEN]{RTC\_CNTL\_DG\_WRAP\_PD\_EN} 使能时，数字调压器在芯片进入 Light-sleep 和 Deep-sleep 时关闭；
            \item 否则，数字调压器常开。
        \end{itemize}
        \item xpd\_peri：
        \begin{itemize}
            \item 当 \hyperref[fielddesc:RTCCNTLDGWRAPPDEN]{RTC\_CNTL\_RTC\_PD\_EN} 使能时，RTC 外设在芯片进入 Light-sleep 和 Deep-sleep 时关闭；
            \item 否则，RTC 外设常开。
        \end{itemize}

        \item xpd\_cpu：
        \begin{itemize}
            \item 当 \hyperref[fielddesc:RTCCNTLDGWRAPPDEN]{RTC\_CNTL\_CPU\_PD\_EN} 使能时，CPU 在芯片进入 Light-sleep 和 Deep-sleep 时关闭；
            \item 否则 CPU 常开。
        \end{itemize}

        \item xpd\_pd\_peri：
        \begin{itemize}
            \item 当 \hyperref[fielddesc:RTCCNTLDGWRAPPDEN]{RTC\_CNTL\_DG\_PERI\_PD\_EN} 使能时，PD 外设在芯片进入 Light-sleep 和 Deep-sleep 时关闭；
            \item 否则，PD 外设常开。
        \end{itemize}

        \item xpd\_dg\_wrap：该开关的控制与 xpd\_dig\_reg 保持一致

        \item xpd\_wireless：
        \begin{itemize}
            \item 当 \hyperref[fielddesc:RTCCNTLDGWRAPPDEN]{RTC\_CNTL\_WIFI\_PD\_EN} 使能时，Wireless circuit 在芯片进入 Light-sleep 和 Deep-sleep 时关闭；
            \item 否则 Wireless circuit 常开。
        \end{itemize}

        \item xpd\_sdio\_reg：见下方第 \ref{sec:lowpowm-flash} 小节。
         \item xpd\_ex\_crystal:
        \begin{itemize}
            \item 当 \hyperref[fielddesc:RTCCNTLXTLFORCEPU]{RTC\_CNTL\_XTL\_FORCE\_PU} 置 1 时，外部主晶振常开；
            \item 否则，外部主晶振在芯片进入 Light-sleep 和 Deep-sleep 时关闭。
        \end{itemize}
        \item xpd\_rc\_oscilator:
        \begin{itemize}
            \item 当 \hyperref[fielddesc:RTCCNTLCK8MFORCEPU]{RTC\_CNTL\_CK8M\_FORCE\_PU} 置 1 时，快速 RC 振荡器常开；
            \item 否则，快速 RC 振荡器在芯片进入 Light-sleep 和 Deep-sleep 时关闭。
        \end{itemize}
        \item RF 电路和 PLL 由内部信号控制，不对用户开放。
        \end{itemize}
\end{itemize}

\subsection{功耗管理单元}
\chipname{} 功耗管理单元可以控制向不同电源域的供电，其主要组成部分包括：

\begin{itemize}
\item RTC 主状态机 (RTC Main State Machine)：产生电源门控、时钟门控和复位信号。
\item 功耗控制器 (Power Controller)：根据 RTC 主状态机产生的电源门控信号，打开或关闭各电源域。
\item 睡眠和唤醒控制器 (Sleep Controller, Wakeup Controller)：向 RTC 主状态机发送睡眠或唤醒请求。
\item 时钟控制器 (Clock Controller)：选择并打开或关闭时钟源。
\item 保护定时器 (Protection Timer)：控制主状态机切换状态的等待时间。
\end{itemize}

在 \chipname{} 的电源管理单元中，睡眠和唤醒控制器向 RTC 主状态机发送睡眠或唤醒请求，RTC 主状态机接着产生电源门控、时钟门控和复位信号。此后，电源控制器和时钟控制器会根据 RTC 主状态机产生的信号，打开或关闭不同的电源域和时钟信号，从而让芯片进入或退出低功耗模式。电源管理单元的主要工作流程可见图 \ref{fig:lpm-rtc-structure}。

\begin{figure}[H]
\centering
\includegraphics[width=0.8\textwidth]{\modulefiles/figures/28-LPM-rtc_structure.pdf}
\caption{电源管理单元的主要工作流程}
\label{fig:lpm-rtc-structure}
\end{figure}

\begin{tiplisting}
\vspace{-3em}
\begin{enumerate}
    \item 完整电源域列表，请见第 \ref{sec:lpm-power-domain} 节。
    \item 有关各唤醒源的具体描述，请见表 \ref{tab:lpm-wakeup-source}。
\end{enumerate}
\end{tiplisting}

\subsection{低功耗时钟}

通常情况下，当 \chipname{} 处于低功耗模式下，芯片的外部晶振 XTAL\_CLK 和 PLL 将被断电以降低功耗，但低功耗时钟仍保持开启，为芯片的不同电源域提供时钟，比如电源管理单元、RTC 外设等，以确保芯片在低功耗模式下的正常工作。

\begin{figure}[H]
\centering
\includegraphics[width=0.8\textwidth]{\modulefiles/figures/28-LPM-rtc_clock.pdf}
\caption{RTC 时钟}
\label{fig:lpm-RTC-clocks}
\end{figure}

\begin{figure}[H]
\centering
\includegraphics[width=0.8\textwidth]{\modulefiles/figures/28-LPM-low_power_clk.pdf}
\caption{Wireless 时钟}
\label{fig:lpm-Digital-low-power-clocks}
\end{figure}


\begin{table}[H]
\caption{低功耗时钟}\label{tab:lpm-clock-type}
\begin{tabular}{| l |L{4cm} | l | l| }
\hline
\rowcolor{lightgray}
\textbf{时钟类型} & \textbf{可选时钟源}  & \textbf{时钟选择寄存器}                                                      & \textbf{作用范围} \\ \hline

\multirow{3}*{RTC 慢速时钟}    & XTAL32K\_CLK &  & 功耗管理单元 \\ \cline{2-2}\cline{4-4}
                              & RC\_FAST\_DIV\_CLK &  &  RTC 定时器 \\ \cline{2-2} \cline{4-4}
 & RC\_SLOW\_CLK（默认）& \multirow{-3}{*}{\hyperref[fielddesc:RTCCNTLANACLKRTCSEL]{RTC\_CNTL\_ANA\_CLK\_RTC\_SEL}} &  ULP 定时器 \\\hline

\multirow{4}*{RTC 快速时钟} & RC\_FAST\_CLK 的 n 分频 &  & ULP 协处理器 \\\cline{4-4}
&（默认） &  & 传感器控制器 \\\cline{2-2} \cline{4-4}
&  &  & RTC 寄存器 \\ \cline{4-4}
 & \multirow{-2}{*}{XTAL\_DIV\_CLK} & \multirow{-4}{*}{\hyperref[fielddesc:RTCCNTLFASTCLKRTCSEL]{RTC\_CNTL\_FAST\_CLK\_RTC\_SEL}} & RTC 慢速内存 \\\hline


\multirow{4}*{低功耗时钟}& XTAL32K\_CLK & \hyperref[regdesc:SYSTEMBTLPCKDIVFRACREG]{SYSTEM\_LPCLK\_SEL\_XTAL32K} & \multirow{4}*{\shortstack{低功耗模式下的数字中\\的无线通信模块 (Wi-Fi/BT)}} \\ \cline{2-3}
& RC\_FAST\_CLK 的 n 分频  & \hyperref[regdesc:SYSTEMBTLPCKDIVFRACREG]{SYSTEM\_LPCLK\_SEL\_8M} &  \\ \cline{2-3}
& RTC\_SLOW\_CLK & \hyperref[fielddesc:SYSTEMLPCLKSELRTCSLOW]{SYSTEM\_LPCLK\_SEL\_RTC\_SLOW} &  \\ \cline{2-3}
 & XTAL\_CLK & \hyperref[regdesc:SYSTEMBTLPCKDIVFRACREG]{SYSTEM\_LPCLK\_SEL\_XTAL} &  \\\hline
\end{tabular}
\end{table}

更多有关时钟的描述，请见章节
\ref{mod:resclk} \textit{\nameref{mod:resclk}}。

\subsection{定时器}

\chipname{} 的低功耗管理使用三个定时器：
\begin{itemize}
\item RTC 定时器
\item ULP 定时器
\item 触摸定时器
\end{itemize}

本章节主要介绍 RTC 定时器。有关 ULP 定时器的内容，请见章节 \ref{mod:ulp} \textit{\nameref{mod:ulp}}。有关触摸定时器的内容，请见章节 \ref{mod:sensor} \textit{\nameref{mod:sensor}}。

RTC 定时器是一个 48 位的可读计数器，可通过配置，使用 RTC 慢速时钟记录以下任一事件发生的时刻。更多详情，请见表 \ref{tab:lpm-Update-RTC-TIMER-VALUE}。

\begin{longtable}[c]{|p{6cm}|p{9cm}|}
\caption{RTC 定时器的触发条件}
\label{tab:lpm-Update-RTC-TIMER-VALUE}
\\ \hline
\rowcolor{lightgray}
\textbf{使能条件} & \textbf{触发条件} \\
\hline
\hyperref[fielddesc:RTCCNTLTIMERXTLOFF]{RTC\_CNTL\_TIMER\_XTL\_OFF} & RTC 主状态机关闭或打开 XTAL\_CLK 晶振时触发。\\
\hline
\multirow{2}{*}{\hyperref[fielddesc:RTCCNTLTIMERSYSSTALL]{RTC\_CNTL\_TIMER\_SYS\_STALL}} & CPU 进入或退出 stall 状态时触发。该设置可保证 SYS\_TIMER 的时间连续性。\\
\hline
\hyperref[fielddesc:RTCCNTLTIMERSYSRST]{RTC\_CNTL\_TIMER\_SYS\_RST} & 系统复位时触发。\\
\hline
\multirow{2}{*}{\hyperref[fielddesc:RTCCNTLRTCTIMEUPDATE]{RTC\_REG\_TIME\_UPDATE}} & 配置寄存器 \hyperref[fielddesc:RTCCNTLRTCTIMEUPDATE]{RTC\_CNTL\_RTC\_TIME\_UPDATE} 时触发。该触发由 CPU 产生（比如用户）。\\
\hline
\end{longtable}


RTC 定时器会在每次触发时更新两组寄存器。其中第一组寄存器记录本次触发的信息，第二组寄存器记录之前触发的信息。这两组寄存器的具体情况见下：

\begin{itemize}
\item 寄存器组 0 用于记录 RTC 定时器在当前触发下的计数值。
\begin{itemize}
\item \hyperref[regdesc:RTCCNTLRTCTIMEHIGH0REG]{RTC\_CNTL\_RTC\_TIME\_HIGH0\_REG}
\item \hyperref[regdesc:RTCCNTLRTCTIMELOW0REG]{RTC\_CNTL\_RTC\_TIME\_LOW0\_REG}
\end{itemize}
\item 寄存器组 1 用于记录 RTC 定时器在上一次触发下的计数值。
\begin{itemize}
\item \hyperref[regdesc:RTCCNTLRTCTIMEHIGH1REG]{RTC\_CNTL\_RTC\_TIME\_HIGH1\_REG}
\item \hyperref[regdesc:RTCCNTLRTCTIMELOW1REG]{RTC\_CNTL\_RTC\_TIME\_LOW1\_REG}
\end{itemize}
\end{itemize}

每次有新的触发，上一次触发时的记录将从寄存器组 0 移至寄存器组 1（寄存器组 1 中之前的记录将被覆盖），而本次触发的记录将存储在寄存器组 0。因此，RTC 定时器最多可同时记录两次触发的值。

值得注意的是，除上电复位外的其余任何复位／睡眠均不会使 RTC 定时器停止或复位。

此外，RTC 定时器还能用作唤醒源。更多详情，请见第 \ref{sec:lpm-wakeup-source} 节。

\subsection{调压器}

\chipname{} 共有三个调压器，负责调节向不同电源域的供电：
\begin{itemize}
\item 数字调压器：负责数字类电源域；
\item 低功耗调压器：负责 RTC 类电源域；
\item flash 调压器：负责数字类和 RTC 类之外的电源域。
\end{itemize}

\begin{tiplisting}
更多有关不同电源域的描述，请见第 \ref{sec:lpm-power-domain} 节。
\end{tiplisting}


\subsubsection{数字调压器}

\chipname{} 的内置数字调压器可以将外部电源电压（通常为 3.3 V）转换为 1.1 V，支持数字类电源域的正常工作，该寄存器主要 xpd\_dig\_reg 信号控制，详见 \ref{fig:lpm-esp32-power-control}。具体结构示意图可见下方图 \ref{fig:lpm-digital-regulator}。

\begin{figure}[H]
\centering
\includegraphics[width=0.5\textwidth]{\modulefiles/figures/28-LPM-REG_DIGITAL.jpg}
\caption{数字调压器}
\label{fig:lpm-digital-regulator}
\end{figure}

\subsubsection{低功耗调压器}

\chipname{} 的内置低功耗调压器可以将外部电源电压（通常为 3.3 V）转换为 1.1 V，支持 RTC 类电源域的正常工作。当管脚 CHIP\_PU 为高电平时，RTC 电路不会掉电。低功耗调压器在芯片进入 Light-sleep 和 Deep-sleep 时可以选择关闭。此时，RTC 电路由一个极低功耗的内置电源供电（该电源无法关闭）。具体结构示意图可见下方图 \ref{fig:lpm-esp32-power-control}。


\begin{figure}[H]
\centering
\includegraphics[width=0.5\textwidth]{\modulefiles/figures/28-LPM-REG_RTC.pdf}
\caption{低功耗调压器}
\label{fig:lpm-low-power-regulator}
\end{figure}


\subsubsection{flash 调压器}\label{sec:lowpowm-flash}

\chipname{} 的内置 flash 调压器可以向数字类和 RTC 类之外的电源域（比如 flash）输出 3.3 V 或 1.8 V 电压，比如 flash，具体结构示意图可见下方图 \ref{fig:lpm-flash-regulator}。

\begin{figure}[H]
\centering
\includegraphics[width=0.7\textwidth]{\modulefiles/figures/28-LPM-REG1P8.pdf}
\caption{flash 调压器}
\label{fig:lpm-flash-regulator}
\end{figure}
\begin{enumerate}
    \item 配置 XPD\_SDIO\_REG 选择调压器供电或外部电源供电：
    \begin{itemize}
        \item 1：由调压器的输出电压供电
        \item 0：由外部电源供电
    \end{itemize}
    \item 配置 SDIO\_TIEH 选择 3.3 V 或 1.8 V 电压
    \begin{itemize}
        \item 1：调压器将管脚 VDD\_SPI 和管脚 VDD3P3\_RTC 短路，电压输出为 3.3 V，即管脚 VDD3P3\_RTC 的电压。
        \item 0：调压器的电压输出为参考电压 VREF，通常为 1.8 V。
    \end{itemize}
\end{enumerate}


上述信号的具体配置如下：
\begin{itemize}
\item XPD\_SDIO\_REG 的配置：
\begin{itemize}
\item 当芯片处于 Active 状态，且\hyperref[fielddesc:RTCCNTLSDIOFORCE]{RTC\_CNTL\_SDIO\_FORCE}  == 0，\hyperref[fielddesc:EFUSEVDDSPIFORCE]{EFUSE\_VDD\_SPI\_FORCE} == 1 时，XPD\_SDIO\_REG 由 \hyperref[fielddesc:EFUSEVDDSPIXPD]{EFUSE\_VDD\_SPI\_XPD} 决定；
\item 当芯片处于 sleep 状态，且\hyperref[fielddesc:RTCCNTLSDIOREGPDEN]{ RTC\_CNTL\_SDIO\_REG\_PD\_EN} == 1 时，XPD\_SDIO\_REG 为 0；
\item 当 \hyperref[fielddesc:RTCCNTLSDIOFORCE]{RTC\_CNTL\_SDIO\_FORCE} == 1 时，XPD\_SDIO\_REG 由 \hyperref[fielddesc:RTCCNTLXPDSDIOREG]{ RTC\_CNTL\_XPD\_SDIO\_REG} 决定。
\end{itemize}
\item SDIO\_TIEH 的配置：
\begin{itemize}
\item 当 \hyperref[fielddesc:RTCCNTLSDIOFORCE]{RTC\_CNTL\_SDIO\_FORCE} == 0 且 \hyperref[fielddesc:EFUSEVDDSPIFORCE]{EFUSE\_VDD\_SPI\_FORCE} == 1 时，SDIO\_TIEH = \hyperref[fielddesc:EFUSEVDDSPITIEH]{EFUSE\_VDD\_SPI\_TIEH} ；
\item 否则，SDIO\_TIEH =  \hyperref[fielddesc:RTCCNTLSDIOTIEH]{RTC\_CNTL\_SDIO\_TIEH}。
\end{itemize}
\end{itemize}

\subsubsection{欠压检测器}

\chipname{} 的欠压检测器可以检查管脚 VDD3P3，VDD3P3\_RTC 和 VDD3P3\_CPU 的电压，在电压快速下落至预设阈值（默认为 2.7 V）以下时发出触发信号，并进行相应处理，从而关闭部分耗电模块（比如 LNA 和 PA 等），为数字模块争取更多时间，用以保存、转移重要数据。

欠压检测器的功耗非常低，在芯片开启时将永远保持开启。
\chipname{} 欠压检测器的具体结构示意图可见下方图 \ref{fig:lpm-brownout-detector}。

\begin{figure}[H]
\centering
\includegraphics[width=0.7\textwidth]{\modulefiles/figures/28-LPM-BROWN_OUT_DET.pdf}
\caption{欠压检测器}
\label{fig:lpm-brownout-detector}
\end{figure}

\hyperref[fielddesc:RTCCNTLRTCBROWNOUTDET]{RTC\_CNTL\_RTC\_BROWN\_OUT\_DET}
可用于指示欠压检测器的输出电平，默认为低电平，可在检测管脚电压下降至阈值以下时跳至高电平。



欠压检测器检测到欠压信号后，有两种处理方法：

\begin{itemize}
    \item mode0：当欠压计数器达到 int comparer 和 rst comparer 中设定的阈值后触发中断，并根据 rst\_sel 的配置选择复位方式。此模式需通过 bod\_mode0\_en 使能。
    \item mode1：直接触发系统复位。
\end{itemize}

详情如下图所示。

\begin{figure}[H]
    \centering
    \includegraphics[width=1.0\textwidth]{\modulefiles/figures/10-LOWPOWM_brownout.pdf}
   \caption{欠压检测器处理方式}
    \label{fig:lpm-Brown-out-detector-mode}
\end{figure}

图中的选择信号均由寄存器控制，对应关系如下：
    \begin{itemize}
        \item bod\_mode0\_en：\hyperref[fielddesc:RTCCNTLBROWNOUTENA]{RTC\_CNTL\_BROWN\_OUT\_ENA}
        \item bod\_mode0\_rst\_en： \hyperref[fielddesc:RTCCNTLBROWNOUTRSTENA]{RTC\_CNTL\_BROWN\_OUT\_RST\_ENA}
        \item bod\_mode0\_en: \hyperref[fielddesc:RTCCNTLBROWNOUTENA]{RTC\_CNTL\_BROWN\_OUT\_ENA}
        \item bod\_mode0\_rst\_sel： \hyperref[fielddesc:RTCCNTLBROWNOUTRSTSEL]{RTC\_CNTL\_BROWN\_OUT\_RST\_SEL} 可用于选择欠压检测器被触发后的复位方式。
            \begin{itemize}
            \item 0：芯片复位
            \item 1：系统复位
            \end{itemize}
        \item bod\_mode1\_sel：\hyperref[fielddesc:RTCCNTLRTCFIBSEL]{RTC\_CNTL\_RTC\_FIB\_SEL} 的第一位。
        \item bod\_mode1\_rst\_en： \hyperref[fielddesc:RTCCNTLBROWNOUTANARSTEN]{RTC\_CNTL\_BROWN\_OUT\_ANA\_RST\_EN}
    \end{itemize}

\section{功耗模式管理}

\subsection{电源域}\label{sec:lpm-power-domain}

\chipname{} 共有三大类共 10 个电源域：
\begin{itemize}
\item RTC 类
    \begin{itemize}
        \item 功耗管理单元
        \item RTC 外设，包括 RTC GPIO、RTC I2C、温度传感器、触摸传感器、RTC ADC 控制器、ULP 协处理器
    \end{itemize}
\item 数字类
    \begin{itemize}
        \item 数字内核
        \item 无线通信 MAC 和 BB
        \item CPU
        \item 部分数字外设，包括 SPI2、GDMA、SHA、RSA、AES、HMAC、DS、Secure Boot、SDIO HOST、USB OTG 等
    \end{itemize}
\item 模拟类
\begin{itemize}
    \item 快速 RC 振荡器 (RC\_FAST\_CLK)
    \item 外部晶振 (XTAL\_CLK)
    \item 锁相环 (PLL)
    \item RF 电路
\end{itemize}
\end{itemize}

\subsection{RTC 状态}\label{sec:lpm-rtc-status}

\chipname{} 共有活跃 (Active)、监测 (Monitor) 和睡眠 (Sleep) 三个主要 RTC 状态，其相互转换过程可见图 \ref{fig:lpm-rtc-states}。

\begin{figure}[H]
\centering
\fbox{\includegraphics[width=0.5\textwidth]{\modulefiles/figures/28-LPM-rtc_states.pdf}}
\caption{RTC 状态}
\label{fig:lpm-rtc-states}
\end{figure}

在不同 RTC 状态下，各电源域的默认开关状态不同，但值得注意的是，用户还可以根据具体需求，强制打开 (FPU) 或关闭 (FPD) 特定电源域。具体可见表 \ref{tab:lpm-rtc-power-domains}。

\begin{table}[H]
\centering
\caption{RTC 状态转换}\label{tab:lpm-rtc-power-domains}
\begin{threeparttable}
\begin{tabular}{|L{3cm}|L{5cm}|l|l|l|}
\hline
\rowcolor{lightgray}
\multicolumn{2}{|c|}{\textbf{电源域}} & \multicolumn{3}{c|}{\textbf{RTC 状态}}  \\\hline
\rowcolor{lightgray}
\textbf{类型} & \textbf{子类型} & \textbf{Active} & \textbf{Monitor} & \textbf{Sleep} \\\hline
\multirow{2}*{RTC\tnote{1，}  \tnote{  2}}                       & 电源管理单元\tnote{3}  & ON & ON  & ON  \\\cline{2-5}
                       & RTC 外设\tnote{4}    & ON & ON  & OFF \\ \hline
\multirow{4}*{数字类}  & CPU\tnote{5}  & ON & OFF* & OFF* \\ \cline{2-5}
 & 无线数字电路\tnote{6} & ON & OFF* & OFF* \\  \cline{2-5}
                       & 数字内核 \tnote{7}& ON & OFF* & OFF* \\ \cline{2-5}
                       & 部分数字外设  & ON & OFF* & OFF* \\ \hline
 \multirow{4}*{模拟类}                       & RC\_FAST\_CLK  & ON & ON  & OFF \\ \cline{2-5}
                       & XTAL\_CLK & ON & OFF & OFF \\ \cline{2-5}
                       & PLL       & ON & OFF & OFF \\ \cline{2-5}
 & RF 电路      & -  & -   & -   \\ \hline
\end{tabular}
\begin{tablenotes}
    \item[*] 可配置。
    \item[1] RTC 慢速内存支持 8 KB SRAM，可用作保留内存或存储 ULP 协处理器指令和数据内存，因此应强制打开。CPU 可通过 PIF 总线访问慢速内存，起始地址为 0x{}5000\_0000。RTC 慢速内存在 Monitor 状态下通常处于 OFF 状态，但这里有一个例外情况，即当 ULP 协处理器工作时，RTC 慢速内存还是处于 ON 的状态。
    \item[2] RTC 快速内存支持 8 KB SRAM，可用作保留内存，因此应强制打开。CPU 可通过 IRAM0/DRAM0 访问快速内存。
    \item[3] \chipname{} 的电源管理单元经过专门设计，一旦芯片上电即处于 ``always-on"（常开）状态，因此无法 FPU 和 FPD。
    \item[4] RTC 外设包括 \ref{mod:ulp} \textit{\nameref{mod:ulp}} 和 \textit{\nameref{mod:sensor}}（比如温度传感器控制器、触摸传感器和 SAR ADC 控制器）。
    \item[5] CPU 在 light-sleep 下也可以单独掉电，但需要 retention DMA 保持恢复 CPU 状态。
    \item[6] 无线通信模块电源域包括 Wi-Fi/BT 各自的 MAC 和 BB（基带）。
    \item[7] 当数字内核关闭时，所有数字类电源域均关闭。值得注意的是，\chipname{} 的 ROM 与 SRAM 也属于数字内核电源域。因此，当数字内核关闭后，用户不可以单独强制打开 ROM 或 SRAM。
\end{tablenotes}
\end{threeparttable}

\end{table}


\subsection{预设功耗模式}
如上文所示，\chipname{} 定义了四种最常见的电源域设置组合，对应 4 种预设功耗模式，可满足用户的常见场景需求，详见表 \ref{tab:lpm-predefined-power-modes}。
\begin{table}[H]
    \centering
    \caption{预设功耗模式}
    \label{tab:lpm-predefined-power-modes}
    \begin{threeparttable}
    \begin{tabular}{|l|c|c|c|c|l|c|c|c|c|c|}
    \hline
        \rowcolor{lightgray}
         & \multicolumn{10}{c|}{\textbf{电源域}} \\\cline{2-11}
        \rowcolor{lightgray}
        \multirow{-2}{*}{\textbf{功耗模式}} & \textbf{PMU} & \textbf{RTC 外设} & \textbf{数字内核} & \textbf{CPU} & \textbf{\makecell{部分数\\字外设}} & \textbf{\makecell{无线数\\字电路}} & \textbf{\makecell{RC\_FAST\\\_CLK}} & \textbf{XTAL\_CLK} & \textbf{PLL} & \textbf{RF 电路} \\\hline
        Active      & ON & ON  & ON  & ON  & ON  & ON   & ON  & ON  & ON  & ON  \\ \hline
        Modem-sleep & ON & ON  & ON  & ON  & ON  & ON\tnote{*} & ON  & ON  & ON  & OFF \\ \hline
        Light-sleep & ON & ON  & ON  & OFF\tnote{*}  & ON\tnote{*}  & OFF\tnote{*} & OFF & OFF & OFF & OFF \\ \hline
        Deep-sleep  & ON & ON\tnote{*}  & OFF  & OFF  & OFF & OFF     & OFF & OFF & OFF & OFF \\ \hline
    \end{tabular}
    \begin{tablenotes}
        \item[*] 可配置
    \end{tablenotes}
    \end{threeparttable}
\end{table}

默认情况下，\chipname{} 系统复位后将进入 Modem-sleep 模式，当有收发包任务时，可配置为 Active 模式。当 CPU 停止工作一段时间后，可以进入 Modem-sleep、Light-sleep 和 Deep-sleep 等低功耗模式。从 Active 到 Deep-sleep 模式，性能递减 \hyperref[other:lpm-predefined-modes-notes]{$^1$}、功耗递减\hyperref[other:lpm-predefined-modes-notes]{$^2$}、唤醒延迟时间递增。此外，这些模式可支持的唤醒源 \hyperref[other:lpm-predefined-modes-notes]{$^3$} 不同。用户根据具体需求，从性能要求、功耗高低、唤醒延迟及可用唤醒源等方面考虑，选择合适的功耗模式。

\begin{tiplisting}
\vspace{-2em}
\label{other:lpm-predefined-modes-notes}
\begin{enumerate}
    \item 更多详情，请见表 \ref{tab:lpm-predefined-power-modes}。
    \item 具体功耗数据可见\href{https://www.espressif.com/sites/default/files/documentation/esp32-s3_datasheet_cn.pdf}{《\chipname{} 技术规格书》}中的功耗特性章节。
    \item 具体可支持的唤醒源，请见第 \ref{sec:lpm-wakeup-source} 节。
\end{enumerate}
\end{tiplisting}


\subsection{唤醒源}\label{sec:lpm-wakeup-source}

\chipname{} 可支持多种唤醒源将 CPU 从不同睡眠模式中唤醒。唤醒源的选择由 \hyperref[fielddesc:RTCCNTLRTCWAKEUPENA]{RTC\_CNTL\_RTC\_WAKEUP\_ENA} 决定，见表 \ref{tab:lpm-wakeup-source}。

\begin{table}[H]
\centering
\caption{唤醒源}
\label{tab:lpm-wakeup-source}
\begin{tabular}[h]{ | r | l | c | c | r | }
\hline
\rowcolor{lightgray}
\textbf{WAKEUP\_ENA} & \textbf{唤醒源\hyperref[wakeup-sources-note10]{$^{10}$}} &  \textbf{Light-sleep} & \textbf{Deep-sleep}  & \textbf{说明}\\
\hline
0x{}1 & EXT0 & Y & Y  & \hyperref[wakeup-sources-note1]{1} \\
\hline
0x{}2 & EXT1 & Y & Y  & \hyperref[wakeup-sources-note2]{2} \\
\hline
0x{}4 & GPIO & Y & Y  & \hyperref[wakeup-sources-note3]{3} \\
\hline
0x{}8 & RTC 定时器 & Y & Y  & -\\
\hline
0x{}20 & Wi-Fi & Y & -  & \hyperref[wakeup-sources-note4]{4} \\
\hline
0x{}40 & UART0 & Y & -  & \hyperref[wakeup-sources-note5]{5} \\
\hline
0x{}80 & UART1 & Y & -  & \hyperref[wakeup-sources-note5]{5} \\
\hline
0x{}100 & TOUCH  Active & Y &  Y &  \hyperref[wakeup-sources-note6]{6}\\
\hline
0x{}200 & ULP-FSM & Y & Y   & \hyperref[wakeup-sources-note7]{7} \\
\hline
0x{}400 & BT & Y & -  & \hyperref[wakeup-sources-note4]{4}\\
\hline
0x{}800 & ULP-RISC-V & Y & Y   &  \\\hline
0x{}1000 & XTAL\_32K & Y & Y   & \hyperref[wakeup-sources-note8]{8} \\
\hline
0x{}2000 & ULP-RISC-V Trap & Y & Y  & \hyperref[wakeup-sources-note9]{9} \\
\hline
0x{}8000 & TOUCH Timeout & Y & Y  & - \\\hline
0x{}c000 & BROWNOUT & Y & Y & - \\\hline
\end{tabular}
\end{table}
\begin{tiplisting}
\vspace{-2em}
\begin{enumerate}
\item \label{wakeup-sources-note1}EXT0 仅可将芯片从 Light-sleep/Deep-sleep 模式中唤醒。当 \hyperref[fielddesc:RTCCNTLEXTWAKEUP0LV]{RTC\_CNTL\_EXT\_WAKEUP0\_LV} 为 1 时将触发管脚高电平，否则触发管脚低电平。用户可通过设置 RTCIO\_EXT\_WAKEUP0\_SEL 选择作为唤醒源的 RTC 管脚。
\item \label{wakeup-sources-note2}EXT1 经过专门设计，可将芯片从任何睡眠模式中唤醒，而且还支持多个管脚的组合。首先，应按照选定作为唤醒源的管脚位图，配置 \hyperref[fielddesc:RTCCNTLEXTWAKEUP1SEL]{RTC\_CNTL\_EXT\_WAKEUP1\_SEL[17:0]}。当 \hyperref[fielddesc:RTCCNTLEXTWAKEUP1LV]{RTC\_CNTL\_EXT\_WAKEUP1\_LV} == 1 时，任何选中管脚之一为高电压则将触发信号唤醒芯片；当 \hyperref[fielddesc:RTCCNTLEXTWAKEUP1LV]{RTC\_CNTL\_REG\_EXT\_WAKEUP1\_LV} == 0 时，任何选中管脚为低电压才能触发信号唤醒芯片。{\bfseries 注意：}用户应确保 EXT1 信号的 hold 时长至少三个 RTC\_SLOW\_CLK 周期，否则信号状态无法被 \hyperref[regdesc:RTCCNTLRTCEXTWAKEUP1STATUSREG]{RTC\_CNTL\_EXT\_WAKEUP1\_STATUS} 记录。
\item \label{wakeup-sources-note3}在 Deep-sleep 模式下，该唤醒源仅有 RTC GPIO 可以作为唤醒源。
\item \label{wakeup-sources-note4}为了通过 Wi-Fi 或 BT 唤醒芯片，芯片将在 Active、Modem-sleep 和 Light-sleep 之间进行切换，CPU 和 RF 模块均将在预设间隔中唤醒，保证 Wi-Fi 和 BT 的正常连接和数据通信。
\item \label{wakeup-sources-note5}当接收到的 RX 脉冲数量超过阈值寄存器中的设置时，即触发唤醒。
\item \label{wakeup-sources-note6}当触摸传感器检测到触摸时，即触发唤醒。
\item \label{wakeup-sources-note7}当超低功耗协处理器配置寄存器  \hyperref[fielddesc:RTCCNTLRTCSWCPUINT]{RTC\_CNTL\_RTC\_SW\_CPU\_INT}，即触发唤醒。
\item \label{wakeup-sources-note8}32 kHz 晶振作为 RTC 慢速时钟时，当 32 kHz 看门狗定时器检测到任何时钟停振，即触发唤醒。
\item \label{wakeup-sources-note9}当超低功耗协处理器进入捕获异常事件时（比如堆栈溢出），即触发唤醒。
\item \label{wakeup-sources-note10}表格中的所有唤醒源均可配置为拒绝睡眠原因，但 UART 除外。
\end{enumerate}
\end{tiplisting}
\subsection{拒绝睡眠}\label{sec:lpm-reject-source}

\chipname{} 提供了硬件拒绝睡眠的机制，防止系统在某些外设仍在工作但是未被 CPU 检测到时进入睡眠，最终导致该外设不能正常工作。

\chipname{} 的所有唤醒源也同时可以作为芯片拒绝睡眠的原因，因此具体的拒绝睡眠源请见表 \ref{tab:lpm-wakeup-source}（UART 除外）。

用户可以配置以下寄存器使能或禁用拒绝睡眠功能。
\begin{itemize}
    \item 配置 \hyperref[fielddesc:RTCCNTLRTCSLEEPREJECTENA]{RTC\_CNTL\_RTC\_SLEEP\_REJECT\_ENA} 整体使能或关闭拒绝休眠功能：
        \begin{itemize}
            \item 进一步配置 \hyperref[fielddesc:RTCCNTLLIGHTSLPREJECTEN]{RTC\_CNTL\_LIGHT\_SLP\_REJECT\_EN}，具体使能拒绝进入 light-sleep；
            \item 进一步配置 \hyperref[fielddesc:RTCCNTLDEEPSLPREJECTEN]{RTC\_CNTL\_DEEP\_SLP\_REJECT\_EN}，具体使能拒绝进入 deep-sleep；
        \end{itemize}
    \item 读取 \hyperref[fielddesc:RTCCNTLREJECTCAUSE]{RTC\_CNTL\_SLP\_REJECT\_CAUSE\_REG} 了解拒绝睡眠的原因。
\end{itemize}

\section{Retention DMA}

\chipname{} 可以选择将 CPU 掉电，进一步降低 light-sleep 模式下的功耗。为了解决 CPU 掉电之后从 light-sleep 唤醒时，程序无法接着从进入睡眠之前的断点运行，\chipname{} 专门新增了 retention 模块。
\chipname{} 的 retention 模块可以在芯片进入 light-sleep 前通过 retention DMA 将 CPU 中的信息备份到 Internal SRAM 的第 2\verb+~+8 块的任意地址中，然后在退出 light-sleep 后将这些信息从 Internal SRAM 中恢复至 CPU，从而保证 CPU 可以接着睡眠之前的断点运行。
\chipname{} 中 retention DMA 的具体特性如下：

\begin{itemize}
    \item Retention DMA 的传输位宽为 128 位，并且只支持 4 字 (Word) 对齐的地址访问。
    \item Retention DMA 的链表经过专门设计，读写操作均共用同一个链表，与通用 DMA 链表的配置方法保持一致：
    \begin{enumerate}
        \item 用户需要在进入睡眠前申请一定的 SRAM 内存空间 \hyperref[dma-note1]{*}，用于存放 CPU 的寄存器信息（428 个字）和配置信息（4 个字），共 432 个字;
        \item 根据内存的申请情况，填写链表的地址和长度信息。具体可参考章节 \ref{mod:dma} \textit{\nameref{mod:dma}}。
    \end{enumerate}

\begin{tiplisting}\label{dma-note1}
\textbf{*}如果申请的 SRAM 空间不够 432 个字，芯片仅能进入普通的 light-sleep，无法进一步关闭 CPU。
\end{tiplisting}
\end{itemize}

此后，用户即可使能寄存器 \hyperref[regdesc:RTCCNTLRETENTIONCTRLREG]{RTC\_CNTL\_RETENTION\_CTRL\_REG} 中的 \hyperref[fielddesc:RTCCNTLRETENTIONEN]{RTC\_CNTL\_RETENTION\_EN} 域，开启 retention 功能，即：
\begin{itemize}
    \item 在系统进入睡眠前，自动启动保留 DMA 进行数据备份；
    \item 在系统唤醒后但未恢复前，自动启动保留 DMA 进行数据还原。
\end{itemize}


\section{RTC Boot}
在 Deep-sleep 下，芯片的 ROM 和 RAM 均将断电，因此在唤醒时 SPI 启动（从 flash 复制数据）所需时间更长。因此，相较于 Light-sleep 和 Modem-sleep 模式，Deep-sleep 模式的唤醒时间要长的多。不过，值得注意的是，在 Deep-sleep 模式下，RTC 快速内存和慢速内存均可以处于上电状态。因此，用户可以将一些代码规模不大（即小于 8 KB 的 ``deep sleep wake stub" ）写入 RTC 快速内存或慢速内存，避免 SPI 启动带来的延迟，从而加速芯片唤醒过程。

\textbf{第一种方法：RTC 慢速内存}
\begin{enumerate}
\item 设置寄存器 \hyperref[fielddesc:RTCCNTLPROCPUSTATVECTORSEL]{RTC\_CNTL\_PROCPU\_STAT\_VECTOR\_SEL} 为 0。
\item 芯片进入睡眠模式。
\item 当 CPU 开启时，复位向量将从地址 0x{}50000000，而非 0x{}40000400 开始复位，整个过程并不需要进行 SPI 启动。RTC 内存中的代码仅需在 C 语言环境中进行部分初始化操作即可。
\end{enumerate}


\textbf{第二种方法：RTC 快速内存}
\begin{enumerate}
\item 设置寄存器 \hyperref[fielddesc:RTCCNTLPROCPUSTATVECTORSEL]{RTC\_CNTL\_PROCPU\_STAT\_VECTOR\_SEL} 为 1。
\item 计算 RTC 快速内存的 CRC 码，并将结果保存在寄存器 \hyperref[regdesc:RTCCNTLRTCSTORE7REG]{RTC\_CNTL\_RTC\_STORE7\_REG[31:0]} 中。
\item 设置寄存器 \hyperref[regdesc:RTCCNTLRTCSTORE6REG]{RTC\_CNTL\_RTC\_STORE6\_REG[31:0]} 为 RTC 快速内存的入口地址。
\item 芯片进入睡眠模式。
\item 当 CPU 开启时，开始进行 ROM 解包和部分初始化工作。此后，再次计算 RTC 快速内存的 CRC 码。如果与寄存器 \hyperref[regdesc:RTCCNTLRTCSTORE7REG]{RTC\_CNTL\_RTC\_STORE7\_REG[31:0]} 中保存的结果一致，则 CPU 跳转至 RTC 快速内存的入口地址。
\end{enumerate}

\chipname{} 的 RTC 启动流程见下方图 \ref{fig:lpm-rtc-boot}：

\begin{figure}[H]
\centering
\includegraphics[width=0.6\textwidth]{\modulefiles/figures/28-LPM-rtc_boot-2.pdf}
\caption{\chipname{} 启动流程图}
\label{fig:lpm-rtc-boot}
\end{figure}

\clearpage
\section{寄存器列表}\label{sec:lowpowm-reg-summ}
\hypertarget{lowpowm-reg-summ}{}

本小节的所有地址均为相对于低功耗管理基地址的地址偏移量（相对地址），具体基地址请见章节 \ref{mod:sysmem} \textit{\nameref{mod:sysmem}} 中的表 \ref{tab:sysmem-base-address}。

请查看章节 \hyperref[glossary-access-types]{\textit{寄存器的访问类型}}，了解“访问”列缩写的含义。

\subfile{\modulefiles/28-LOWPOWM-reg-summ__CN}

\clearpage
\section{寄存器}

本小节的所有地址均为相对于低功耗管理基地址的地址偏移量（相对地址），具体基地址请见章节 \ref{mod:sysmem} \textit{\nameref{mod:sysmem}} 中的表 \ref{tab:sysmem-base-address}。
\subfile{\modulefiles/28-LOWPOWM-reg__CN}


\end{document}
